\subsection{Simulation Environment Implementation}

In order to facilitate the systematic design and testing of antenna
alignment algorithms before hardware deployment, a simulation environment was
implemented using the Python programming language. This simulation offers a
reproducible approximation of the wireless transmission characteristics,
antenna radiation pattern, and mechanical movement capabilities of the
proposed system.

The simulation environment aims to mimic the iterative process of
communication between an adaptive algorithm and a directional antenna
system. It features a step-by-step interface that can be utilized by both
reinforcement learning algorithms and traditional antenna alignment
approaches.

\subsubsection{Structure of Environment and State Representation}

The environment itself operates on the principles of discrete-time Markov
decision process (MDP). Every time step involves updating the orientation of the
antenna and calculating RSSI for it.

Antenna orientation is characterized as:
\begin{itemize}
    \item Azimuth angle $\theta \in [0^\circ, 360^\circ)$ in equal angular steps 
    \item Elevation (tilt) angle $\phi \in [60^\circ, 120^\circ]$ in $5^\circ$ steps
\end{itemize}

Alongside absolute antenna orientation, the reinforcement learner is provided with an
additional discrete value that describes short term changes in RSSI, namely whether it is rising, steady or declining. 

This leads to the following state description formula:
\begin{equation}
    s = (\theta_i, \phi_j, r_k)
\end{equation}
where $\theta_i$ is an index of a certain azimuth angle step, $\phi_j$ is an index of elevation angle step and $r_k$ indicates RSSI trend.

\subsubsection{RSSI Calculation and Propagation Model}

RSSI values for each orientation of the antenna are calculated using an
idealized but physically motivated propagation model. The value of RSSI for
a particular orientation is determined by path loss, gain and random noise.

Path loss is simulated using a log-distance path loss model:
\begin{equation}
PL(d) = PL(d_0) + 10n \log_{10}\left(\frac{d}{d_0}\right)
\end{equation}
where $d$ denotes distance between the transmitter and the receiver, $d_0$ is a
reference distance, and $n$ is path loss exponent \cite{ref3}. In the present simulations,
distance $d$ is held constant to simulate a stationary transmitter.

RSSI is thus defined as:
\begin{equation}
RSSI = P_t + G_t(\theta, \phi) + G_r - PL(d) + \eta
\end{equation}
where $P_t$ is the transmitting power, $G_t(\theta, \phi)$ is the directional gain
of the receiving antenna as a function of its orientation, and $G_r$ is the gain
of the transmitter antenna.

\subsubsection{Modeling Noise in RSSI Measurements}

For a more realistic simulation of RSSI fluctuations experienced in a Wi-Fi
connection, additive noise is added to RSSI measurements. The noise variable
$\eta$ is assumed to be zero-mean Gaussian:
\begin{equation}
\eta \sim \mathcal{N}(0, \sigma^2)
\end{equation}

Furthermore, moving average filtering is used to simulate the smoothing effect of
packet-level averaging done in the embedded Wi-Fi interface implementation.

\subsubsection{Action Space and Environmental Dynamics}

The environmental state includes support for discrete actions corresponding to small
movements of the antenna. Each movement consists of some change in both azimuth and
tilt angle, constrained by the physical nature of the problem.

At each time step after choosing an action, the environment does the following:
\begin{itemize}
    \item Performs a change in antenna orientation.
    \item Simulates an environment delay for the change to take effect.
    \item Obtains a measurement of RSSI according to the current antenna orientation.
    \item Calculates the reward function based on RSSI increase.
\end{itemize}

This process is used in all experiments, ensuring that the dynamics of the
environment are identical for all algorithms considered.

\subsubsection{Evaluation of Algorithm in the Simulation}

The simulation framework provides a common basis for evaluation of all the alignment
algorithms that have been discussed in this study. In particular, the performance of the reinforcement
learning algorithm, as well as the baseline approaches, is tested on a set of episodes with the same
initial conditions and RSSI noise samples.

Metrics used to evaluate algorithm performance include:
\begin{itemize}
    \item Time required for reaching optimal alignment
    \item Highest reached RSSI level
    \item Number of antenna moves
    \item Resistance to RSSI noise
\end{itemize}
